% Kompilacja: pdflatex main && bibtex main && pdflatex main && pdflatex main
% Dane liczbowe (definicje.tex), tabele (tabele/*.csv) i rysunki (rysunki/) generuje
% analiza_odpowiedzi_naukowa.py; tekst (main.tex, sekcja-*.tex) jest edytowany ręcznie.
\documentclass[11pt]{article}
\usepackage[utf8]{inputenc}
\usepackage[T1]{fontenc}
\usepackage[polish]{babel}
\usepackage[margin=2.5cm]{geometry}
\usepackage{booktabs,graphicx,xcolor,amsmath}
\usepackage[hidelinks]{hyperref}
\usepackage{doi}
% BIBLIOGRAFIA:
\usepackage[numbers,sort&compress]{natbib}
\bibliographystyle{plainnat}

\extrafloats{100}
\newcommand{\fillin}[1]{\textcolor{red}{[#1]}}
% Wygenerowano automatycznie przez analiza_odpowiedzi_naukowa.py -- plik jest NADPISYWANY, nie edytuj ręcznie.
% Wstaw w preambule: % Wygenerowano automatycznie przez analiza_odpowiedzi_naukowa.py -- plik jest NADPISYWANY, nie edytuj ręcznie.
% Wstaw w preambule: % Wygenerowano automatycznie przez analiza_odpowiedzi_naukowa.py -- plik jest NADPISYWANY, nie edytuj ręcznie.
% Wstaw w preambule: \input{definicje.tex}
%
% Liczby w notacji polskiej: przecinek dziesiętny {,} (działa w tekście i w trybie matematycznym),
% odstęp tysięcy \, (od 10 000), minus \ensuremath{-}, przedziały [a; b].
% Makra wartości p (\SignP, \WilcoxonP, \P...) zawierają relację i są do trybu matematycznego: $p \SignP$.
%
% Tabele: \tabelacsv{<kolumny tabular>}{<nazwa pliku bez .csv>} wczytuje \TabDir<nazwa>.csv
% (csvsimple-l3; separator ';'; pierwszy wiersz = nagłówek; komórki w składni LaTeX; komórki ze średnikiem
% są ujęte w klamry; prefiks \midrule w pierwszej komórce = linia oddzielająca grupy wierszy).
% Przypis pod tabelą: \uwagatabeli{...}. Katalogi \TabDir i \FigDir można zdefiniować przed \input.
\RequirePackage{booktabs,csvsimple-l3}
\providecommand{\TabDir}{tabele/}
\providecommand{\FigDir}{rysunki/}
\providecommand{\fillin}[1]{\textbf{[#1]}}
\providecommand{\tabelacsv}[2]{\csvloop{file={\TabDir #2.csv}, separator=semicolon, head,
  column names detection=false, after head=\begin{tabular}{#1}\toprule\csvlinetotablerow\\\midrule,
  command=\csvlinetotablerow, late after line=\\, late after last line=\\\bottomrule\end{tabular}}}
\providecommand{\uwagatabeli}[1]{\par\vspace{2pt}\parbox{\linewidth}{\footnotesize #1}}
%
% Stałe liczbowe
\newcommand{\NResponses}{60\,962}
\newcommand{\NEmployees}{187}
\newcommand{\NItems}{164}
\newcommand{\NModules}{9}
\newcommand{\NSequences}{15\,783}
\newcommand{\NRejected}{0}
\newcommand{\DateStart}{kwietnia 2023}
\newcommand{\DateEnd}{czerwca 2026}
\newcommand{\AccOverall}{86{,}5}
\newcommand{\AccFirst}{82{,}5}
\newcommand{\AttemptsMedian}{3}
\newcommand{\AttemptsMax}{19}
\newcommand{\SeqRepeatedPct}{78{,}6}
\newcommand{\SeqMinThreePct}{60{,}6}
\newcommand{\GapMedian}{36}
\newcommand{\GapQOne}{18}
\newcommand{\GapQThree}{90}
\newcommand{\DailyMedian}{3}
\newcommand{\RespPerEmpMedian}{218}
\newcommand{\RespPerEmpQOne}{67}
\newcommand{\RespPerEmpQThree}{592}
\newcommand{\SpanMonthsMedian}{18}
\newcommand{\ActiveLastQuarterPct}{56{,}1}
\newcommand{\WeekdayPct}{100{,}0}
\newcommand{\ItemsPerEmpMedian}{90}
\newcommand{\ActiveDaysMedian}{84}
\newcommand{\OptionsMedian}{3}
\newcommand{\NExcluded}{212}
\newcommand{\NExcludedItems}{2}
\newcommand{\NExcludedItemResp}{348}
\newcommand{\ExcludedModules}{mnk-tst}
\newcommand{\ReaskAfterCorrect}{74{,}0}
\newcommand{\ReaskAfterIncorrect}{74{,}7}
\newcommand{\GapAfterCorrect}{36}
\newcommand{\GapAfterIncorrect}{38}
\newcommand{\AccAttOne}{82{,}5}
\newcommand{\AccAttOneCI}{[81{,}0; 83{,}9]}
\newcommand{\AccAttMax}{91{,}8}
\newcommand{\AccAttMaxCI}{[88{,}5; 94{,}4]}
\newcommand{\AttMax}{10}
\newcommand{\CohortK}{5}
\newcommand{\CohortN}{5187}
\newcommand{\CohortAccOne}{81{,}8}
\newcommand{\CohortAccK}{89{,}1}
\newcommand{\CohortAccOneCI}{[79{,}7; 83{,}8]}
\newcommand{\CohortAccKCI}{[87{,}3; 90{,}6]}
\newcommand{\PowerLawB}{0{,}29}
\newcommand{\PowerLawRsq}{0{,}98}
\newcommand{\ORMOneAttempt}{1{,}38}
\newcommand{\ORMOneAttemptCI}{[1{,}32; 1{,}45]}
\newcommand{\PMOneAttempt}{< 0{,}001}
\newcommand{\ORMTwoAttempt}{1{,}33}
\newcommand{\ORMTwoAttemptCI}{[1{,}24; 1{,}42]}
\newcommand{\PMTwoAttempt}{< 0{,}001}
\newcommand{\ORMTwoPrev}{3{,}85}
\newcommand{\ORMTwoPrevCI}{[3{,}49; 4{,}26]}
\newcommand{\PMTwoPrev}{< 0{,}001}
\newcommand{\ORMTwoGap}{0{,}88}
\newcommand{\ORMTwoGapCI}{[0{,}86; 0{,}90]}
\newcommand{\PMTwoGap}{< 0{,}001}
\newcommand{\ORMTwoTenure}{0{,}87}
\newcommand{\ORMTwoTenureCI}{[0{,}77; 1{,}00]}
\newcommand{\PMTwoTenure}{= 0{,}044}
\newcommand{\ORMThreeTenure}{0{,}88}
\newcommand{\ORMThreeTenureCI}{[0{,}76; 1{,}01]}
\newcommand{\PMThreeTenure}{= 0{,}061}
\newcommand{\MOneDropped}{282}
\newcommand{\ICCEmployee}{0{,}09}
\newcommand{\ICCItem}{0{,}24}
\newcommand{\SDEmployee}{0{,}69}
\newcommand{\SDItem}{1{,}10}
\newcommand{\RsqMarginal}{0{,}021}
\newcommand{\RsqConditional}{0{,}35}
\newcommand{\ORRange}{1{,}29--1{,}50}
\newcommand{\ORSensMain}{1{,}38}
\newcommand{\ORSensMainCI}{[1{,}32; 1{,}45]}
\newcommand{\ORSensCohort}{1{,}36}
\newcommand{\ORSensCohortCI}{[1{,}27; 1{,}46]}
\newcommand{\ORSensTrunc}{1{,}38}
\newcommand{\ORSensTruncCI}{[1{,}32; 1{,}44]}
\newcommand{\ORSensLenThree}{1{,}30}
\newcommand{\ORSensLenThreeCI}{[1{,}21; 1{,}40]}
\newcommand{\ORSensLenFive}{1{,}35}
\newcommand{\ORSensLenFiveCI}{[1{,}26; 1{,}45]}
\newcommand{\ORSensLenEight}{1{,}40}
\newcommand{\ORSensLenEightCI}{[1{,}31; 1{,}50]}
\newcommand{\ORSensActive}{1{,}38}
\newcommand{\ORSensActiveCI}{[1{,}32; 1{,}45]}
\newcommand{\ORSensTenure}{1{,}50}
\newcommand{\ORSensTenureCI}{[1{,}42; 1{,}59]}
\newcommand{\ORSensGee}{1{,}29}
\newcommand{\ORSensGeeCI}{[1{,}24; 1{,}35]}
\newcommand{\ORSensMixed}{1{,}37}
\newcommand{\ORSensMixedCI}{[1{,}34; 1{,}41]}
\newcommand{\RetainPct}{92{,}6}
\newcommand{\RetainCI}{[91{,}7; 93{,}3]}
\newcommand{\ForgetPct}{7{,}4}
\newcommand{\RecoverPct}{58{,}5}
\newcommand{\RecoverCI}{[55{,}7; 61{,}3]}
\newcommand{\PersistPct}{41{,}5}
\newcommand{\SameDistractorPct}{71{,}6}
\newcommand{\SameDistractorCI}{[69{,}4; 73{,}8]}
\newcommand{\SameDistractorN}{2546}
\newcommand{\SameDistractorExpected}{47{,}1}
\newcommand{\SameDistractorExpectedPop}{59{,}6}
\newcommand{\RetShortCorrect}{95{,}2}
\newcommand{\RetLongCorrect}{88{,}5}
\newcommand{\RetShortIncorrect}{61{,}0}
\newcommand{\RetLongIncorrect}{58{,}4}
\newcommand{\EmpN}{149}
\newcommand{\EmpImproved}{102}
\newcommand{\EmpWorsened}{20}
\newcommand{\EmpUnchanged}{27}
\newcommand{\EmpMeanChange}{3{,}8}
\newcommand{\EmpMeanChangeCI}{[2{,}4; 5{,}2]}
\newcommand{\SignP}{< 0{,}001}
\newcommand{\WilcoxonP}{< 0{,}001}
\newcommand{\RankBiserial}{0{,}72}
\newcommand{\EmpSharePct}{79{,}7}
\newcommand{\ItemDifficultyMedian}{85{,}3}
\newcommand{\TopDistractorMedian}{71{,}8}
\newcommand{\NDistractors}{349}
\newcommand{\NNonFunctional}{193}
\newcommand{\NonFunctionalPct}{55{,}3}
\newcommand{\ItemsModerate}{102}
\newcommand{\ItemsGoodDisc}{49}
\newcommand{\ItemsDiscTested}{161}
\newcommand{\ItemsImproving}{44}
\newcommand{\ItemsWorsening}{1}
\newcommand{\ItemsTestable}{164}
\newcommand{\ItemDiscMedian}{0{,}22}
\newcommand{\ItemsLowDisc}{70}
\newcommand{\ItemsEasy}{58}
\newcommand{\ItemsHard}{4}
\newcommand{\TopItems}{10}
\newcommand{\DemoN}{187}
\newcommand{\DemoMatched}{151}
\newcommand{\DemoMatchedPct}{80{,}7}
\newcommand{\DemoUnmatched}{36}
\newcommand{\DemoWomen}{121}
\newcommand{\DemoWomenPct}{80{,}1}
\newcommand{\DemoMen}{30}
\newcommand{\DemoMenPct}{19{,}9}
\newcommand{\DemoSexMissing}{36}
\newcommand{\DemoAgeN}{150}
\newcommand{\DemoAgeMissing}{37}
\newcommand{\AgeMedian}{47}
\newcommand{\AgeQOne}{41}
\newcommand{\AgeQThree}{56}
\newcommand{\AgeMin}{22}
\newcommand{\AgeMax}{67}
\newcommand{\AgeMean}{47{,}2}
\newcommand{\AgeSD}{10{,}6}
\newcommand{\AgeModalGroup}{50--59}
\newcommand{\AgeModalGroupPct}{34{,}0}
\newcommand{\BootReps}{2000}
\newcommand{\Alpha}{0{,}05}
\newcommand{\MinItemN}{30}
\newcommand{\Seed}{2026}
\newcommand{\MaxAttemptPooled}{10}
\newcommand{\VerPython}{3.14.4}
\newcommand{\VerPandas}{3.0.3}
\newcommand{\VerNumpy}{2.5.0}
\newcommand{\VerScipy}{1.18.0}
\newcommand{\VerStatsmodels}{0.14.6}
\newcommand{\VerMatplotlib}{3.11.0}

%
% Liczby w notacji polskiej: przecinek dziesiętny {,} (działa w tekście i w trybie matematycznym),
% odstęp tysięcy \, (od 10 000), minus \ensuremath{-}, przedziały [a; b].
% Makra wartości p (\SignP, \WilcoxonP, \P...) zawierają relację i są do trybu matematycznego: $p \SignP$.
%
% Tabele: \tabelacsv{<kolumny tabular>}{<nazwa pliku bez .csv>} wczytuje \TabDir<nazwa>.csv
% (csvsimple-l3; separator ';'; pierwszy wiersz = nagłówek; komórki w składni LaTeX; komórki ze średnikiem
% są ujęte w klamry; prefiks \midrule w pierwszej komórce = linia oddzielająca grupy wierszy).
% Przypis pod tabelą: \uwagatabeli{...}. Katalogi \TabDir i \FigDir można zdefiniować przed \input.
\RequirePackage{booktabs,csvsimple-l3}
\providecommand{\TabDir}{tabele/}
\providecommand{\FigDir}{rysunki/}
\providecommand{\fillin}[1]{\textbf{[#1]}}
\providecommand{\tabelacsv}[2]{\csvloop{file={\TabDir #2.csv}, separator=semicolon, head,
  column names detection=false, after head=\begin{tabular}{#1}\toprule\csvlinetotablerow\\\midrule,
  command=\csvlinetotablerow, late after line=\\, late after last line=\\\bottomrule\end{tabular}}}
\providecommand{\uwagatabeli}[1]{\par\vspace{2pt}\parbox{\linewidth}{\footnotesize #1}}
%
% Stałe liczbowe
\newcommand{\NResponses}{60\,962}
\newcommand{\NEmployees}{187}
\newcommand{\NItems}{164}
\newcommand{\NModules}{9}
\newcommand{\NSequences}{15\,783}
\newcommand{\NRejected}{0}
\newcommand{\DateStart}{kwietnia 2023}
\newcommand{\DateEnd}{czerwca 2026}
\newcommand{\AccOverall}{86{,}5}
\newcommand{\AccFirst}{82{,}5}
\newcommand{\AttemptsMedian}{3}
\newcommand{\AttemptsMax}{19}
\newcommand{\SeqRepeatedPct}{78{,}6}
\newcommand{\SeqMinThreePct}{60{,}6}
\newcommand{\GapMedian}{36}
\newcommand{\GapQOne}{18}
\newcommand{\GapQThree}{90}
\newcommand{\DailyMedian}{3}
\newcommand{\RespPerEmpMedian}{218}
\newcommand{\RespPerEmpQOne}{67}
\newcommand{\RespPerEmpQThree}{592}
\newcommand{\SpanMonthsMedian}{18}
\newcommand{\ActiveLastQuarterPct}{56{,}1}
\newcommand{\WeekdayPct}{100{,}0}
\newcommand{\ItemsPerEmpMedian}{90}
\newcommand{\ActiveDaysMedian}{84}
\newcommand{\OptionsMedian}{3}
\newcommand{\NExcluded}{212}
\newcommand{\NExcludedItems}{2}
\newcommand{\NExcludedItemResp}{348}
\newcommand{\ExcludedModules}{mnk-tst}
\newcommand{\ReaskAfterCorrect}{74{,}0}
\newcommand{\ReaskAfterIncorrect}{74{,}7}
\newcommand{\GapAfterCorrect}{36}
\newcommand{\GapAfterIncorrect}{38}
\newcommand{\AccAttOne}{82{,}5}
\newcommand{\AccAttOneCI}{[81{,}0; 83{,}9]}
\newcommand{\AccAttMax}{91{,}8}
\newcommand{\AccAttMaxCI}{[88{,}5; 94{,}4]}
\newcommand{\AttMax}{10}
\newcommand{\CohortK}{5}
\newcommand{\CohortN}{5187}
\newcommand{\CohortAccOne}{81{,}8}
\newcommand{\CohortAccK}{89{,}1}
\newcommand{\CohortAccOneCI}{[79{,}7; 83{,}8]}
\newcommand{\CohortAccKCI}{[87{,}3; 90{,}6]}
\newcommand{\PowerLawB}{0{,}29}
\newcommand{\PowerLawRsq}{0{,}98}
\newcommand{\ORMOneAttempt}{1{,}38}
\newcommand{\ORMOneAttemptCI}{[1{,}32; 1{,}45]}
\newcommand{\PMOneAttempt}{< 0{,}001}
\newcommand{\ORMTwoAttempt}{1{,}33}
\newcommand{\ORMTwoAttemptCI}{[1{,}24; 1{,}42]}
\newcommand{\PMTwoAttempt}{< 0{,}001}
\newcommand{\ORMTwoPrev}{3{,}85}
\newcommand{\ORMTwoPrevCI}{[3{,}49; 4{,}26]}
\newcommand{\PMTwoPrev}{< 0{,}001}
\newcommand{\ORMTwoGap}{0{,}88}
\newcommand{\ORMTwoGapCI}{[0{,}86; 0{,}90]}
\newcommand{\PMTwoGap}{< 0{,}001}
\newcommand{\ORMTwoTenure}{0{,}87}
\newcommand{\ORMTwoTenureCI}{[0{,}77; 1{,}00]}
\newcommand{\PMTwoTenure}{= 0{,}044}
\newcommand{\ORMThreeTenure}{0{,}88}
\newcommand{\ORMThreeTenureCI}{[0{,}76; 1{,}01]}
\newcommand{\PMThreeTenure}{= 0{,}061}
\newcommand{\MOneDropped}{282}
\newcommand{\ICCEmployee}{0{,}09}
\newcommand{\ICCItem}{0{,}24}
\newcommand{\SDEmployee}{0{,}69}
\newcommand{\SDItem}{1{,}10}
\newcommand{\RsqMarginal}{0{,}021}
\newcommand{\RsqConditional}{0{,}35}
\newcommand{\ORRange}{1{,}29--1{,}50}
\newcommand{\ORSensMain}{1{,}38}
\newcommand{\ORSensMainCI}{[1{,}32; 1{,}45]}
\newcommand{\ORSensCohort}{1{,}36}
\newcommand{\ORSensCohortCI}{[1{,}27; 1{,}46]}
\newcommand{\ORSensTrunc}{1{,}38}
\newcommand{\ORSensTruncCI}{[1{,}32; 1{,}44]}
\newcommand{\ORSensLenThree}{1{,}30}
\newcommand{\ORSensLenThreeCI}{[1{,}21; 1{,}40]}
\newcommand{\ORSensLenFive}{1{,}35}
\newcommand{\ORSensLenFiveCI}{[1{,}26; 1{,}45]}
\newcommand{\ORSensLenEight}{1{,}40}
\newcommand{\ORSensLenEightCI}{[1{,}31; 1{,}50]}
\newcommand{\ORSensActive}{1{,}38}
\newcommand{\ORSensActiveCI}{[1{,}32; 1{,}45]}
\newcommand{\ORSensTenure}{1{,}50}
\newcommand{\ORSensTenureCI}{[1{,}42; 1{,}59]}
\newcommand{\ORSensGee}{1{,}29}
\newcommand{\ORSensGeeCI}{[1{,}24; 1{,}35]}
\newcommand{\ORSensMixed}{1{,}37}
\newcommand{\ORSensMixedCI}{[1{,}34; 1{,}41]}
\newcommand{\RetainPct}{92{,}6}
\newcommand{\RetainCI}{[91{,}7; 93{,}3]}
\newcommand{\ForgetPct}{7{,}4}
\newcommand{\RecoverPct}{58{,}5}
\newcommand{\RecoverCI}{[55{,}7; 61{,}3]}
\newcommand{\PersistPct}{41{,}5}
\newcommand{\SameDistractorPct}{71{,}6}
\newcommand{\SameDistractorCI}{[69{,}4; 73{,}8]}
\newcommand{\SameDistractorN}{2546}
\newcommand{\SameDistractorExpected}{47{,}1}
\newcommand{\SameDistractorExpectedPop}{59{,}6}
\newcommand{\RetShortCorrect}{95{,}2}
\newcommand{\RetLongCorrect}{88{,}5}
\newcommand{\RetShortIncorrect}{61{,}0}
\newcommand{\RetLongIncorrect}{58{,}4}
\newcommand{\EmpN}{149}
\newcommand{\EmpImproved}{102}
\newcommand{\EmpWorsened}{20}
\newcommand{\EmpUnchanged}{27}
\newcommand{\EmpMeanChange}{3{,}8}
\newcommand{\EmpMeanChangeCI}{[2{,}4; 5{,}2]}
\newcommand{\SignP}{< 0{,}001}
\newcommand{\WilcoxonP}{< 0{,}001}
\newcommand{\RankBiserial}{0{,}72}
\newcommand{\EmpSharePct}{79{,}7}
\newcommand{\ItemDifficultyMedian}{85{,}3}
\newcommand{\TopDistractorMedian}{71{,}8}
\newcommand{\NDistractors}{349}
\newcommand{\NNonFunctional}{193}
\newcommand{\NonFunctionalPct}{55{,}3}
\newcommand{\ItemsModerate}{102}
\newcommand{\ItemsGoodDisc}{49}
\newcommand{\ItemsDiscTested}{161}
\newcommand{\ItemsImproving}{44}
\newcommand{\ItemsWorsening}{1}
\newcommand{\ItemsTestable}{164}
\newcommand{\ItemDiscMedian}{0{,}22}
\newcommand{\ItemsLowDisc}{70}
\newcommand{\ItemsEasy}{58}
\newcommand{\ItemsHard}{4}
\newcommand{\TopItems}{10}
\newcommand{\DemoN}{187}
\newcommand{\DemoMatched}{151}
\newcommand{\DemoMatchedPct}{80{,}7}
\newcommand{\DemoUnmatched}{36}
\newcommand{\DemoWomen}{121}
\newcommand{\DemoWomenPct}{80{,}1}
\newcommand{\DemoMen}{30}
\newcommand{\DemoMenPct}{19{,}9}
\newcommand{\DemoSexMissing}{36}
\newcommand{\DemoAgeN}{150}
\newcommand{\DemoAgeMissing}{37}
\newcommand{\AgeMedian}{47}
\newcommand{\AgeQOne}{41}
\newcommand{\AgeQThree}{56}
\newcommand{\AgeMin}{22}
\newcommand{\AgeMax}{67}
\newcommand{\AgeMean}{47{,}2}
\newcommand{\AgeSD}{10{,}6}
\newcommand{\AgeModalGroup}{50--59}
\newcommand{\AgeModalGroupPct}{34{,}0}
\newcommand{\BootReps}{2000}
\newcommand{\Alpha}{0{,}05}
\newcommand{\MinItemN}{30}
\newcommand{\Seed}{2026}
\newcommand{\MaxAttemptPooled}{10}
\newcommand{\VerPython}{3.14.4}
\newcommand{\VerPandas}{3.0.3}
\newcommand{\VerNumpy}{2.5.0}
\newcommand{\VerScipy}{1.18.0}
\newcommand{\VerStatsmodels}{0.14.6}
\newcommand{\VerMatplotlib}{3.11.0}

%
% Liczby w notacji polskiej: przecinek dziesiętny {,} (działa w tekście i w trybie matematycznym),
% odstęp tysięcy \, (od 10 000), minus \ensuremath{-}, przedziały [a; b].
% Makra wartości p (\SignP, \WilcoxonP, \P...) zawierają relację i są do trybu matematycznego: $p \SignP$.
%
% Tabele: \tabelacsv{<kolumny tabular>}{<nazwa pliku bez .csv>} wczytuje \TabDir<nazwa>.csv
% (csvsimple-l3; separator ';'; pierwszy wiersz = nagłówek; komórki w składni LaTeX; komórki ze średnikiem
% są ujęte w klamry; prefiks \midrule w pierwszej komórce = linia oddzielająca grupy wierszy).
% Przypis pod tabelą: \uwagatabeli{...}. Katalogi \TabDir i \FigDir można zdefiniować przed \input.
\RequirePackage{booktabs,csvsimple-l3}
\providecommand{\TabDir}{tabele/}
\providecommand{\FigDir}{rysunki/}
\providecommand{\fillin}[1]{\textbf{[#1]}}
\providecommand{\tabelacsv}[2]{\csvloop{file={\TabDir #2.csv}, separator=semicolon, head,
  column names detection=false, after head=\begin{tabular}{#1}\toprule\csvlinetotablerow\\\midrule,
  command=\csvlinetotablerow, late after line=\\, late after last line=\\\bottomrule\end{tabular}}}
\providecommand{\uwagatabeli}[1]{\par\vspace{2pt}\parbox{\linewidth}{\footnotesize #1}}
%
% Stałe liczbowe
\newcommand{\NResponses}{60\,962}
\newcommand{\NEmployees}{187}
\newcommand{\NItems}{164}
\newcommand{\NModules}{9}
\newcommand{\NSequences}{15\,783}
\newcommand{\NRejected}{0}
\newcommand{\DateStart}{kwietnia 2023}
\newcommand{\DateEnd}{czerwca 2026}
\newcommand{\AccOverall}{86{,}5}
\newcommand{\AccFirst}{82{,}5}
\newcommand{\AttemptsMedian}{3}
\newcommand{\AttemptsMax}{19}
\newcommand{\SeqRepeatedPct}{78{,}6}
\newcommand{\SeqMinThreePct}{60{,}6}
\newcommand{\GapMedian}{36}
\newcommand{\GapQOne}{18}
\newcommand{\GapQThree}{90}
\newcommand{\DailyMedian}{3}
\newcommand{\RespPerEmpMedian}{218}
\newcommand{\RespPerEmpQOne}{67}
\newcommand{\RespPerEmpQThree}{592}
\newcommand{\SpanMonthsMedian}{18}
\newcommand{\ActiveLastQuarterPct}{56{,}1}
\newcommand{\WeekdayPct}{100{,}0}
\newcommand{\ItemsPerEmpMedian}{90}
\newcommand{\ActiveDaysMedian}{84}
\newcommand{\OptionsMedian}{3}
\newcommand{\NExcluded}{212}
\newcommand{\NExcludedItems}{2}
\newcommand{\NExcludedItemResp}{348}
\newcommand{\ExcludedModules}{mnk-tst}
\newcommand{\ReaskAfterCorrect}{74{,}0}
\newcommand{\ReaskAfterIncorrect}{74{,}7}
\newcommand{\GapAfterCorrect}{36}
\newcommand{\GapAfterIncorrect}{38}
\newcommand{\AccAttOne}{82{,}5}
\newcommand{\AccAttOneCI}{[81{,}0; 83{,}9]}
\newcommand{\AccAttMax}{91{,}8}
\newcommand{\AccAttMaxCI}{[88{,}5; 94{,}4]}
\newcommand{\AttMax}{10}
\newcommand{\CohortK}{5}
\newcommand{\CohortN}{5187}
\newcommand{\CohortAccOne}{81{,}8}
\newcommand{\CohortAccK}{89{,}1}
\newcommand{\CohortAccOneCI}{[79{,}7; 83{,}8]}
\newcommand{\CohortAccKCI}{[87{,}3; 90{,}6]}
\newcommand{\PowerLawB}{0{,}29}
\newcommand{\PowerLawRsq}{0{,}98}
\newcommand{\ORMOneAttempt}{1{,}38}
\newcommand{\ORMOneAttemptCI}{[1{,}32; 1{,}45]}
\newcommand{\PMOneAttempt}{< 0{,}001}
\newcommand{\ORMTwoAttempt}{1{,}33}
\newcommand{\ORMTwoAttemptCI}{[1{,}24; 1{,}42]}
\newcommand{\PMTwoAttempt}{< 0{,}001}
\newcommand{\ORMTwoPrev}{3{,}85}
\newcommand{\ORMTwoPrevCI}{[3{,}49; 4{,}26]}
\newcommand{\PMTwoPrev}{< 0{,}001}
\newcommand{\ORMTwoGap}{0{,}88}
\newcommand{\ORMTwoGapCI}{[0{,}86; 0{,}90]}
\newcommand{\PMTwoGap}{< 0{,}001}
\newcommand{\ORMTwoTenure}{0{,}87}
\newcommand{\ORMTwoTenureCI}{[0{,}77; 1{,}00]}
\newcommand{\PMTwoTenure}{= 0{,}044}
\newcommand{\ORMThreeTenure}{0{,}88}
\newcommand{\ORMThreeTenureCI}{[0{,}76; 1{,}01]}
\newcommand{\PMThreeTenure}{= 0{,}061}
\newcommand{\MOneDropped}{282}
\newcommand{\ICCEmployee}{0{,}09}
\newcommand{\ICCItem}{0{,}24}
\newcommand{\SDEmployee}{0{,}69}
\newcommand{\SDItem}{1{,}10}
\newcommand{\RsqMarginal}{0{,}021}
\newcommand{\RsqConditional}{0{,}35}
\newcommand{\ORRange}{1{,}29--1{,}50}
\newcommand{\ORSensMain}{1{,}38}
\newcommand{\ORSensMainCI}{[1{,}32; 1{,}45]}
\newcommand{\ORSensCohort}{1{,}36}
\newcommand{\ORSensCohortCI}{[1{,}27; 1{,}46]}
\newcommand{\ORSensTrunc}{1{,}38}
\newcommand{\ORSensTruncCI}{[1{,}32; 1{,}44]}
\newcommand{\ORSensLenThree}{1{,}30}
\newcommand{\ORSensLenThreeCI}{[1{,}21; 1{,}40]}
\newcommand{\ORSensLenFive}{1{,}35}
\newcommand{\ORSensLenFiveCI}{[1{,}26; 1{,}45]}
\newcommand{\ORSensLenEight}{1{,}40}
\newcommand{\ORSensLenEightCI}{[1{,}31; 1{,}50]}
\newcommand{\ORSensActive}{1{,}38}
\newcommand{\ORSensActiveCI}{[1{,}32; 1{,}45]}
\newcommand{\ORSensTenure}{1{,}50}
\newcommand{\ORSensTenureCI}{[1{,}42; 1{,}59]}
\newcommand{\ORSensGee}{1{,}29}
\newcommand{\ORSensGeeCI}{[1{,}24; 1{,}35]}
\newcommand{\ORSensMixed}{1{,}37}
\newcommand{\ORSensMixedCI}{[1{,}34; 1{,}41]}
\newcommand{\RetainPct}{92{,}6}
\newcommand{\RetainCI}{[91{,}7; 93{,}3]}
\newcommand{\ForgetPct}{7{,}4}
\newcommand{\RecoverPct}{58{,}5}
\newcommand{\RecoverCI}{[55{,}7; 61{,}3]}
\newcommand{\PersistPct}{41{,}5}
\newcommand{\SameDistractorPct}{71{,}6}
\newcommand{\SameDistractorCI}{[69{,}4; 73{,}8]}
\newcommand{\SameDistractorN}{2546}
\newcommand{\SameDistractorExpected}{47{,}1}
\newcommand{\SameDistractorExpectedPop}{59{,}6}
\newcommand{\RetShortCorrect}{95{,}2}
\newcommand{\RetLongCorrect}{88{,}5}
\newcommand{\RetShortIncorrect}{61{,}0}
\newcommand{\RetLongIncorrect}{58{,}4}
\newcommand{\EmpN}{149}
\newcommand{\EmpImproved}{102}
\newcommand{\EmpWorsened}{20}
\newcommand{\EmpUnchanged}{27}
\newcommand{\EmpMeanChange}{3{,}8}
\newcommand{\EmpMeanChangeCI}{[2{,}4; 5{,}2]}
\newcommand{\SignP}{< 0{,}001}
\newcommand{\WilcoxonP}{< 0{,}001}
\newcommand{\RankBiserial}{0{,}72}
\newcommand{\EmpSharePct}{79{,}7}
\newcommand{\ItemDifficultyMedian}{85{,}3}
\newcommand{\TopDistractorMedian}{71{,}8}
\newcommand{\NDistractors}{349}
\newcommand{\NNonFunctional}{193}
\newcommand{\NonFunctionalPct}{55{,}3}
\newcommand{\ItemsModerate}{102}
\newcommand{\ItemsGoodDisc}{49}
\newcommand{\ItemsDiscTested}{161}
\newcommand{\ItemsImproving}{44}
\newcommand{\ItemsWorsening}{1}
\newcommand{\ItemsTestable}{164}
\newcommand{\ItemDiscMedian}{0{,}22}
\newcommand{\ItemsLowDisc}{70}
\newcommand{\ItemsEasy}{58}
\newcommand{\ItemsHard}{4}
\newcommand{\TopItems}{10}
\newcommand{\DemoN}{187}
\newcommand{\DemoMatched}{151}
\newcommand{\DemoMatchedPct}{80{,}7}
\newcommand{\DemoUnmatched}{36}
\newcommand{\DemoWomen}{121}
\newcommand{\DemoWomenPct}{80{,}1}
\newcommand{\DemoMen}{30}
\newcommand{\DemoMenPct}{19{,}9}
\newcommand{\DemoSexMissing}{36}
\newcommand{\DemoAgeN}{150}
\newcommand{\DemoAgeMissing}{37}
\newcommand{\AgeMedian}{47}
\newcommand{\AgeQOne}{41}
\newcommand{\AgeQThree}{56}
\newcommand{\AgeMin}{22}
\newcommand{\AgeMax}{67}
\newcommand{\AgeMean}{47{,}2}
\newcommand{\AgeSD}{10{,}6}
\newcommand{\AgeModalGroup}{50--59}
\newcommand{\AgeModalGroupPct}{34{,}0}
\newcommand{\BootReps}{2000}
\newcommand{\Alpha}{0{,}05}
\newcommand{\MinItemN}{30}
\newcommand{\Seed}{2026}
\newcommand{\MaxAttemptPooled}{10}
\newcommand{\VerPython}{3.14.4}
\newcommand{\VerPandas}{3.0.3}
\newcommand{\VerNumpy}{2.5.0}
\newcommand{\VerScipy}{1.18.0}
\newcommand{\VerStatsmodels}{0.14.6}
\newcommand{\VerMatplotlib}{3.11.0}


\usepackage{authblk}
\renewcommand\Authand{ i }\renewcommand\Authands{ i }
\author[1,*]{Bartosz Lewandowski \small{(ORCID:0009-0004-1158-0693)}}
\author[2]{Monika Szerszeń \small{(ORCID: 0009-0009-6625-3979)}}
\author[3]{Magdalena Dowgalec \small{(ORCID:0009-0005-4469-2832)}}
\author[4]{Ewaryst Tkacz \small{(ORCID:0000-0003-2580-7954)}}
\author[5]{Maciej J. Wróbel \small{(ORCID:0000-0003-1865-2122)}}

\affil[1,4]{Department of Clinical Engineering, Academy of Silesia in Katowice}
\affil[3]{Powiślański Bank Spółdzieczy w Kwidzynie}
\affil[4]{Atende SA}
\affil[5]{Department of Otolaryngology and Oncological Laryngology Collegium Medicum NCU}

% Autor korespondencyjny
\affil[*]{Autor korespondencyjny: dr Bartosz Lewandowski, bartosz.lewandowski@akademiaslaska.pl}

\title{Przyrost wiedzy bez transferu: skuteczność mikrouczenia z powtórkami rozłożonymi
w czasie w szkoleniach pracowniczych}

\newcommand{\Women}{xx}


\date{}
\begin{document}
\maketitle


\begin{abstract}
\noindent\textbf{Wprowadzenie.} Przypominanie rozłożone w czasie poprawia zapamiętywanie w warunkach
laboratoryjnych i edukacyjnych, jednak wyniki wieloletnich programów szkoleń w miejscu pracy opisano
w niewielu badaniach.
\textbf{Metody.} Przeanalizowano \NResponses{} odpowiedzi udzielonych przez \NEmployees{} pracowników
(wśród osób o znanych danych \DemoWomenPct\% kobiet, mediana wieku \AgeMedian{} lat)
na \NItems{} pytań wielokrotnego wyboru w \NModules{} modułach programu mikrouczenia opartego na
quizach (od \DateStart{} do \DateEnd). Każde pytanie było ponownie zadawane temu samemu pracownikowi
w stałym, nieadaptacyjnym odstępie (mediana \GapMedian{} dni). Poprawność odpowiedzi modelowano
regresją logistyczną z efektami stałymi pracownika i pytania; oszacowania opisowe podano
z przedziałami ufności z bootstrapu klastrowego.
\textbf{Wyniki.} Odsetek poprawnych odpowiedzi wzrósł z \AccAttOne\% w pierwszej próbie do
\AccAttMax\% w próbie \AttMax{}. i kolejnych. Każde podwojenie liczby prób zwiększało szanse
poprawnej odpowiedzi \ORMOneAttempt-krotnie (95\% PU \ORMOneAttemptCI; zakres w dziesięciu
specyfikacjach analiz wrażliwości \ORRange). Dłuższy odstęp od poprzedniej próby wiązał się z niższą
poprawnością (OR \ORMTwoGap{} na podwojenie odstępu). Powtarzane błędy koncentrowały się na tym samym
dystraktorze (\SameDistractorPct\% wobec oczekiwanych \SameDistractorExpectedPop\%). Poprawność
odpowiedzi na nowo napotkane pytania nie rosła wraz ze stażem (OR \ORMThreeTenure{} na rok). % SPRAWDŹ
\textbf{Wnioski.} Powtarzanemu przypominaniu w ramach rutynowych szkoleń pracowniczych towarzyszyła
poprawa ograniczona do powtarzanych pytań oraz zapominanie między powtórkami. Skuteczność programu
mogłyby zwiększyć adaptacyjny harmonogram powtórek, informacja zwrotna odnosząca się do utrwalonych
błędnych przekonań oraz modyfikacja pytań o niskiej dyskryminacji i dystraktorów niefunkcjonalnych.

\medskip\noindent\textbf{Słowa kluczowe:} powtórki rozłożone w czasie; praktyka przypominania;
mikrouczenie; szkolenia pracownicze; analityka uczenia się; analiza pytań testowych
\end{abstract}

% =============================================================================
% Wprowadzenie. Plik edytowany ręcznie -- skrypt analiza_odpowiedzi_naukowa.py go NIE nadpisuje.
% Liczby: definicje.tex; tabele: \TabDir (CSV); rysunki: \FigDir; literatura: references.bib.
% =============================================================================
\section{Wprowadzenie}

Przypominanie sobie informacji z pamięci wzmacnia ich późniejsze utrwalenie silniej niż ponowne
ich przeczytanie (efekt testowania), a nauka rozłożona w czasie jest zapamiętywana lepiej niż nauka
skumulowana (efekt rozłożenia w czasie) \citep{roediger2006,cepeda2006,cepeda2008}. W randomizowanych
badaniach z udziałem lekarzy w trakcie kształcenia tzw. edukacja rozłożona w czasie (ang.
\emph{spaced education}), w której krótkie pytania wielokrotnego wyboru przesyła się cyklicznie
przez tygodnie lub miesiące, poprawiała przyswajanie i utrwalanie wiedzy
\citep{kerfoot2007urology,kerfoot2007mededuc,kerfoot2009}, natomiast przeglądy systematyczne wskazują na małe
lub umiarkowane korzyści u pracowników ochrony zdrowia \citep{phillips2019,martinengo2024}.

W szkoleniach organizacyjnych stosuje się krótkie formy \emph{mikrouczenia} (ang.
\emph{microlearning}) oparte na quizach, jednak dostępne dowody pochodzą głównie z badań na małych
próbach lub z tzw. szarej literatury \citep{taylor2022}. Ewaluacje obowiązkowych szkoleń organizacyjnych, np.\
z zakresu świadomości bezpieczeństwa, często wykazują efekty niewielkie albo zanikające z upływem czasu
\citep{reinheimer2020,lain2022,ho2025}. Dane z dzienników (logów) takich platform pozwalają
natomiast obserwować uczenie się na poziomie pojedynczych odpowiedzi, podobnie jak w modelach ucznia
rozwijanych dla inteligentnych systemów tutorskich i aplikacji do nauki języków
\citep{corbett1994,pavlik2009,settles2016,tabibian2019,koedinger2023}.

Badanie przeprowadzono w banku, w którym pracownicy odpowiadali na krótkie pytania wielokrotnego
wyboru dotyczące wymogów regulacyjnych, produktów bankowych, umiejętności miękkich oraz, w przypadku
kadry kierowniczej, zarządzania, wyświetlane przez aplikację na komputerze w trakcie dnia pracy.
Trzyletnie logi odpowiedzi z tego programu posłużyły do weryfikacji następujących pytań
badawczych:
(PB1) czy odsetek poprawnych odpowiedzi rośnie wraz z kolejnymi przypomnieniami tego samego pytania
przez tego samego pracownika;
(PB2) jak poprawność zależy od wyniku poprzedniej próby i od czasu, który od niej upłynął;
(PB3) czy powtarzane błędy są trwałe, co wskazywałoby raczej na błędne przekonania niż na
zgadywanie;
(PB4) czy poprawa przenosi się na nowo napotykane pytania; oraz
(PB5) jakie są właściwości psychometryczne pytań?

% =============================================================================
% Materiał i metody. Plik edytowany ręcznie -- skrypt analiza_odpowiedzi_naukowa.py go NIE nadpisuje.
% Liczby: definicje.tex; tabele: \TabDir (CSV); rysunki: \FigDir; literatura: references.bib.
% =============================================================================
\section{Materiał i metody}

\subsection{Środowisko i program szkoleniowy}
Program prowadzono w banku. Jego treść stanowiły pytania ogólne pogrupowane w cztery kategorie:
wymogi regulacyjne, produkty bankowe, umiejętności miękkie oraz moduł dla kadry kierowniczej.
Każde pytanie było pytaniem wielokrotnego wyboru z trzema do pięciu opcjami odpowiedzi, z których
dokładnie jedna była poprawna, i zawierało wyjaśnienie poprawnej odpowiedzi. Aplikacja szkoleniowa
uruchamiała się automatycznie na ekranach komputerów pracowników i wyświetlała od jednego do trzech
pytań dziennie w godzinach 10:00--14:00; po każdej odpowiedzi pokazywano poprawną odpowiedź wraz
z wyjaśnieniem. To samo pytanie mogło zostać ponownie zadane temu samemu pracownikowi nie wcześniej
niż siedem dni po poprzedniej odpowiedzi, a harmonogram powtórek nie zależał od tego, czy odpowiedź
była poprawna. Na czas przeglądu i aktualizacji pytań program był zawieszany.
O udziale w programie decydował Departament Digitalizacji i Procesów banku, który był również autorem
pytań. Pytania prezentowano w modułach tematycznych (mediana \OptionsMedian{} zaobserwowanych opcji
odpowiedzi na pytanie). W analizowanych danych pracownicy odpowiadali na medianę \DailyMedian{}
pytań w dniu aktywności, \WeekdayPct\% odpowiedzi udzielono w dni robocze, a to samo pytanie było
ponownie zadawane temu samemu pracownikowi po medianie \GapMedian{} dni
(IQR \GapQOne--\GapQThree).

\subsection{Dane}
Dane zostały zorganizoane w firmie listy zawierającej  moduł, pytanie, wybraną opcję,
kod odpowiadającego, znacznik czasu oraz poprawność odpowiedzi (flaga true/false). Wykluczono  \NExcludedItems{} pytania, w których zarejestrowano
tylko jedną opcję odpowiedzi (\NExcludedItemResp{} odpowiedzi), ponieważ nie niosą one informacji
o błędach. Zbiór analityczny obejmował \NResponses{} odpowiedzi (tab.~\ref{tab:zbior_danych}).
Dane osób odpowowiadających zostały przekazane w formie metadanych, zawierające informacje o płci panelisty oraz przedziale wiekowym (co 10 lat).

\fillin{Podstawa prawna przetwarzania danych spseudonimizowanych (wiek, płeć) oraz zgoda komisji
bioetycznej lub zwolnienie z obowiązku jej uzyskania, jeśli wymaga tego czasopismo.}

\subsection{Definicje}
\emph{Sekwencja} to uporządkowana chronologicznie seria prób jednego pracownika dotyczących jednego
pytania; \emph{numer próby} to pozycja odpowiedzi w jej sekwencji (przy jednakowym czasie o kolejności
decydował identyfikator rekordu). \emph{Poprawność} to odsetek poprawnych odpowiedzi. \emph{Odstęp}
to czas w dniach od poprzedniej próby w tej samej sekwencji. \emph{Staż} to czas od pierwszej
odpowiedzi pracownika w programie. Zmianę na poziomie pracownika zdefiniowano jako różnicę odsetka
poprawnych odpowiedzi między próbą~3 a próbą~1, obliczoną na podstawie własnych sekwencji pracownika
liczących co najmniej trzy próby.

\subsection{Analiza statystyczna}
Odsetki poprawnych odpowiedzi w statystykach opisowych podano z 95-procentowymi przedziałami ufności (PU) z bootstrapu
klastrowego, w którym losowano pracowników ze zwracaniem (\BootReps{} replikacji); procedura ta
uwzględnia zależność odpowiedzi udzielanych przez tę samą osobę. Ponieważ osoby kontynuujące do kolejnych prób mogą różnić się od
tych, które przestają, zagregowane krzywe uczenia mogą być zniekształcone przez ubywanie obserwacji
\citep{murray2013,goutte2018}; dlatego krzywe uczenia obliczono dodatkowo dla zbalansowanej kohorty
sekwencji liczących co najmniej \CohortK{} prób. Zależność odsetka błędów od numeru próby opisano
funkcją potęgową \citep{newell1981}. Ponieważ harmonogram powtórek nie był adaptacyjny, sprawdzono,
czy ponowne zadanie pytania zależało od wyniku poprzedniej próby.

Poprawność odpowiedzi modelowano regresją logistyczną z efektami stałymi pracownika i pytania oraz
błędami standardowymi klastrowanymi względem pracowników, tak aby efekt numeru próby był szacowany
w obrębie osób i w obrębie pytań. Model~1 zawierał numer próby w skali log$_2$, dzięki czemu iloraz
szans (OR) odnosi się do podwojenia liczby prób. Model~2, dopasowany do prób powtórnych, uwzględniał
dodatkowo wynik poprzedniej próby, odstęp od poprzedniej próby (log$_2$ dni) oraz staż, analogicznie
do modeli \emph{performance factors analysis} i \emph{half-life regression} stosowanych do logów
uczących się \citep{pavlik2009,settles2016}. Model~3, dopasowany wyłącznie do pierwszych prób,
weryfikował, czy poprawność odpowiedzi na nowo napotkane pytania zmieniała się wraz ze stażem (PB4).
W analizach wrażliwości wykorzystano zbalansowaną kohortę, obcięcie do 10 prób, warstwy według długości
sekwencji (kontrola obciążenia typu \emph{pattern-mixture}), wykluczenie pracowników o niskiej
aktywności, korektę o staż, uogólnione równania estymujące (GEE) z wymienną macierzą korelacji
roboczej \citep{liang1986} oraz model mieszany ze skrzyżowanymi losowymi wyrazami wolnymi pracowników
i pytań \citep{baayen2008}, na podstawie którego wyznaczono współczynniki korelacji wewnątrzklasowej
(ICC) w skali latentnej oraz brzegowy i warunkowy współczynnik $R^2$ \citep{nakagawa2013}.

Przejścia między kolejnymi próbami zestawiono w tabeli kontyngencji. Dla dwóch następujących po sobie błędów porównano odsetek ponownych wyborów
tego samego dystraktora z odsetkiem oczekiwanym przy losowym wyborze spośród zaobserwowanych
dystraktorów oraz przy wyborze proporcjonalnym do popularności każdego dystraktora w danym pytaniu.
Zmianę na poziomie pracowników testowano testem znaków i testem rangowanych znaków Wilcoxona
(korelacja rangowo-dwuseryjna jako miara wielkości efektu). Pytania opisano wskaźnikami klasycznej
teorii testu obliczonymi na pierwszych próbach: trudnością (odsetek poprawnych odpowiedzi),
dyskryminacją (korelacja pytanie--reszta, tj.\ z poprawnością pracownika w pozostałych pytaniach) oraz
funkcjonowaniem dystraktorów; dystraktory wybierane przez mniej niż 5\% odpowiadających uznano za
niefunkcjonalne \citep{tarrant2009,gierl2017}. Trendy poprawności w kolejnych próbach dla
poszczególnych pytań testowano regresją logistyczną z błędami standardowymi klastrowanymi względem
pracowników, kontrolując odsetek fałszywych odkryć (FDR) na poziomie \Alpha{} \citep{benjamini1995}.
Analizy wykonano w języku Python \VerPython{} z użyciem bibliotek pandas \VerPandas, NumPy \VerNumpy,
SciPy \VerScipy, statsmodels \VerStatsmodels{} i Matplotlib \VerMatplotlib{}.
%  (ziarno generatoraliczb losowych \Seed).

% =============================================================================
% Wyniki. Plik edytowany ręcznie -- skrypt analiza_odpowiedzi_naukowa.py go NIE nadpisuje.
% Liczby: definicje.tex; tabele: \tabelacsv{<kolumny>}{<plik>} wczytuje \TabDir<plik>.csv;
% rysunki: \FigDir. Zdania interpretujące kierunek wyniku oznaczono komentarzem "% SPRAWDŹ" --
% zweryfikuj je po ponownym uruchomieniu skryptu na nowych danych lub z innymi parametrami.
% =============================================================================
\section{Wyniki}

\subsection{Próba i zaangażowanie}
Zbiór danych obejmował \NResponses{} odpowiedzi \NEmployees{} pracowników na \NItems{} pytań
w \NModules{} modułach, zebranych od \DateStart{} do \DateEnd{} (tab.~\ref{tab:zbior_danych}).
Mediana liczby odpowiedzi na pracownika wynosiła \RespPerEmpMedian{} (IQR
\RespPerEmpQOne--\RespPerEmpQThree), udzielonych na \ItemsPerEmpMedian{} różnych pytań w ciągu
\ActiveDaysMedian{} dni aktywności; mediana czasu uczestnictwa w programie wynosiła
\SpanMonthsMedian{} miesięcy, a \ActiveLastQuarterPct\% pracowników było aktywnych w ostatnich
90 dniach obserwacji. Aktywność zmieniała się w czasie i występowały okresy bez odpowiedzi
(rys.~\ref{fig:aktywnosc}); odpowiadają one zawieszeniom programu na czas przeglądu i aktualizacji
pytań. Ogólny odsetek poprawnych odpowiedzi wynosił \AccOverall\%, a w pierwszych próbach \AccFirst\%.

Spośród \DemoN{} pracowników objętych analizą dane o wieku i płci były dostępne dla \DemoMatched{}
(\DemoMatchedPct\%; tab.~\ref{tab:demografia}). Wśród osób o znanej płci kobiety stanowiły
\DemoWomenPct\% ($n = \DemoWomen$), a mężczyźni \DemoMenPct\% ($n = \DemoMen$). Mediana wieku wynosiła
\AgeMedian{} lat (IQR \AgeQOne--\AgeQThree; zakres \AgeMin--\AgeMax), a najliczniejszą grupę stanowiły
osoby w wieku \AgeModalGroup{} lat (\AgeModalGroupPct\% osób o znanym wieku; rys.~\ref{fig:demografia}).
% SPRAWDŹ: przedział modalny i proporcje płci po zmianie danych

Harmonogram powtórek nie zależał od wyniku poprzedniej próby: prawdopodobieństwo kolejnej próby
wynosiło \ReaskAfterCorrect\% po odpowiedzi poprawnej i \ReaskAfterIncorrect\% po błędnej, natomiast mediana
odstępu do kolejnej próby wynosiła odpowiednio \GapAfterCorrect{} i \GapAfterIncorrect{} dni. Kontynuacja
sekwencji nie była więc selekcjonowana według wyników. % SPRAWDŹ

\begin{table}[htbp]
\centering
\caption{Charakterystyka zbioru danych.}
\label{tab:zbior_danych} 
\small
\tabelacsv{lr}{zbior_danych}
\uwagatabeli{IQR: rozstęp międzykwartylowy. Sekwencja to seria prób jednego pracownika dotyczących
jednego pytania. Wykluczone moduły testowe: \ExcludedModules{} (\NExcluded{} odpowiedzi). Wykluczone
pytania z jedną zaobserwowaną opcją odpowiedzi: \NExcludedItems{} (\NExcludedItemResp{} odpowiedzi).}
\end{table}

\begin{table}[htbp]
\centering
\caption{Charakterystyka demograficzna pracowników objętych analizą.}
\label{tab:demografia}
\small
\tabelacsv{lr}{demografia}
\uwagatabeli{Odsetki płci liczono wśród osób o znanej płci, a odsetki przedziałów wieku wśród osób
o znanym wieku ($n = \DemoAgeN$). Brak danych dotyczy głównie pracowników nieobecnych w zestawieniu
demograficznym (\DemoUnmatched). IQR: rozstęp międzykwartylowy; SD: odchylenie standardowe.}
\end{table}

\begin{figure*}[tbp]
\centering
\includegraphics{\FigDir rys_aktywnosc}
\caption{Miesięczna liczba odpowiedzi (u góry) i aktywnych pracowników (u dołu) w okresie obserwacji.}
\label{fig:aktywnosc}
\end{figure*}

\begin{figure}[tbp]
\centering
\includegraphics{\FigDir rys_demografia}
\caption{Pracownicy objęci analizą według wieku (przedziały 10-letnie) i płci. Liczby nad słupkami:
łączna liczba pracowników w przedziale. Pominięto osoby bez danych o wieku ($n = \DemoAgeMissing$).}
\label{fig:demografia}
\end{figure}

\subsection{Uczenie się w kolejnych próbach (PB1)}
Odsetek poprawnych odpowiedzi rósł wraz z numerem próby: z \AccAttOne\%
\AccAttOneCI{} w pierwszej próbie do \AccAttMax\% \AccAttMaxCI{} w próbie \AttMax{}. i kolejnych
(rys.~\ref{fig:krzywa_uczenia}, tab.~\ref{tab:krzywa_uczenia}). Wzrost był podobny w zbalansowanej
kohorcie \CohortN{} sekwencji liczących co najmniej \CohortK{} prób (\CohortAccOne\% \CohortAccOneCI{}
w próbie~1 wobec \CohortAccK\% \CohortAccKCI{} w próbie~\CohortK), nie wynika więc z selektywnej
kontynuacji przez pracowników osiągających lepsze wyniki. Odsetek błędów był funkcją potęgową numeru
próby (wykładnik \PowerLawB, $R^2 = \PowerLawRsq$), zgodnie z potęgowym prawem praktyki
\citep{newell1981}.

W modelu~1 każde podwojenie liczby prób wiązało się z \ORMOneAttempt-krotnie wyższymi szansami
poprawnej odpowiedzi (95\% PU \ORMOneAttemptCI; tab.~\ref{tab:regresja}). Pytania odpowiadały za
większą część wariancji niż pracownicy (ICC w skali latentnej \ICCItem{} wobec \ICCEmployee; brzegowy
$R^2 = \RsqMarginal$, warunkowy $R^2 = \RsqConditional$).

\begin{figure}[tbp]
\centering
\includegraphics{\FigDir rys_krzywa_uczenia}
\caption{Odsetek poprawnych odpowiedzi według numeru próby w sekwencjach pracownik--pytanie. Koła:
wszystkie sekwencje dostępne w danej próbie (ich liczba maleje wraz z numerem próby). Kwadraty:
zbalansowana kohorta, tj.\ te same \CohortN{} sekwencje, które osiągnęły co najmniej \CohortK{} prób,
obserwowane wyłącznie w próbach 1--\CohortK{}, dzięki czemu każdy punkt opiera się na identycznych
sekwencjach. Słupki błędów: 95\% PU z bootstrapu klastrowego po pracownikach (\BootReps{} replikacji).
Linia przerywana: dopasowanie funkcji potęgowej do odsetka błędów. Próby \MaxAttemptPooled{}.
i dalsze połączono.}
\label{fig:krzywa_uczenia}
\end{figure}

\begin{table*}[htbp]
\centering
\caption{Odsetek poprawnych odpowiedzi według numeru próby w sekwencjach pracownik--pytanie.}
\label{tab:krzywa_uczenia}
\small
\setlength{\tabcolsep}{3.5pt}
\tabelacsv{lrrrlrrl}{krzywa_uczenia}
\uwagatabeli{Kohorta: sekwencje liczące $\geq$\CohortK{} prób, co utrzymuje stały zbiór sekwencji
w próbach 1--\CohortK{} i eliminuje obciążenie przeżywalności. 95\% PU z bootstrapu klastrowego,
w którym losowano pracowników (\BootReps{} replikacji).}
\end{table*}

\begin{table*}[htbp]
\centering
\caption{Regresja logistyczna poprawności odpowiedzi: ilorazy szans [95\% PU].}
\label{tab:regresja}
\small
\setlength{\tabcolsep}{3.5pt}
\tabelacsv{llll}{regresja}
\uwagatabeli{Wszystkie modele zawierają efekty stałe pracownika i pytania; błędy standardowe
klastrowane względem pracowników. Pracownicy i pytania bez zmienności poprawności są pomijani przez
warunkową funkcję wiarygodności (model~1: \MOneDropped{} odpowiedzi). Model~3 weryfikuje, czy
poprawność odpowiedzi na nowo napotkane pytania rośnie wraz z czasem uczestnictwa w programie.}
\end{table*}

\subsection{Wynik poprzedniej próby, odstęp i zapominanie (PB2)}
Po uwzględnieniu numeru próby poprawna poprzednia odpowiedź wiązała się z wyższymi szansami poprawnej
odpowiedzi (model~2: OR \ORMTwoPrev, 95\% PU \ORMTwoPrevCI), natomiast dłuższy odstęp od poprzedniej
próby wiązał się z niższymi szansami (OR \ORMTwoGap{} na podwojenie odstępu, 95\% PU \ORMTwoGapCI; tab.~\ref{tab:regresja}).
Wśród odpowiedzi poprzednio poprawnych odsetek poprawnych zmniejszał się z \RetShortCorrect\%, gdy
pytanie powtórzono w ciągu 14 dni, do \RetLongCorrect\% po 240 dniach lub dłużej; odpowiedzi
poprzednio błędne były korygowane odpowiednio w \RetShortIncorrect\% i \RetLongIncorrect\% przypadków
(rys.~\ref{fig:odstep}, tab.~\ref{tab:retencja_wg_odstepu}). Wzorzec ten jest zgodny z zapominaniem
między powtórkami \citep{cepeda2006,cepeda2008}. % SPRAWDŹ

\begin{figure}[tbp]
\centering
\includegraphics{\FigDir rys_odstep}
\caption{Odsetek poprawnych odpowiedzi według odstępu od poprzedniej próby tego samego pytania,
osobno dla odpowiedzi poprzednio poprawnych (koła) i poprzednio błędnych (kwadraty). Słupki błędów:
95\% PU z bootstrapu klastrowego po pracownikach.}
\label{fig:odstep}
\end{figure}

\begin{table}[htbp]
\centering
\caption{Odsetek poprawnych odpowiedzi według odstępu od poprzedniej próby tego samego pytania,
dla odpowiedzi poprzednio poprawnych (po lewej) i błędnych (po prawej).}
\label{tab:retencja_wg_odstepu}
\small
\tabelacsv{lrlrl}{retencja_wg_odstepu}
\uwagatabeli{95\% PU z bootstrapu klastrowego, w którym losowano pracowników (\BootReps{} replikacji).}
\end{table}

\subsection{Przejścia i utrwalone błędy (PB3)}
Spośród odpowiedzi poprzednio poprawnych \RetainPct\% \RetainCI{} było poprawnych również w kolejnej
próbie (\ForgetPct\% zapomniano); spośród odpowiedzi poprzednio błędnych \RecoverPct\% \RecoverCI{}
skorygowano, a \PersistPct\% pozostało błędnych (tab.~\ref{tab:przejscia}). Gdy błąd się powtarzał
(\SameDistractorN{} przypadków), pracownicy wybierali ten sam dystraktor w \SameDistractorPct\%
\SameDistractorCI{} przypadków, wobec \SameDistractorExpected\% oczekiwanych przy wyborze losowym
i \SameDistractorExpectedPop\% oczekiwanych przy uwzględnieniu popularności poszczególnych
dystraktorów, co wskazuje na trwałe błędne przekonania, a nie na zgadywanie. % SPRAWDŹ

\begin{table}[htbp]
\centering
\caption{Przejścia między kolejnymi próbami tego samego pytania.}
\label{tab:przejscia}
\small
\tabelacsv{lrrlr}{przejscia}
\uwagatabeli{Po dwóch kolejnych błędach ten sam dystraktor wybrano w \SameDistractorPct\%
\SameDistractorCI{} przypadków ($n = \SameDistractorN$); wartość oczekiwana przy losowym wyborze spośród
zaobserwowanych dystraktorów wynosiła \SameDistractorExpected\%, a po uwzględnieniu ogólnej popularności
każdego dystraktora w danym pytaniu wynosiła ona \SameDistractorExpectedPop\%. 95\% PU z bootstrapu klastrowego,
w którym losowano pracowników (\BootReps{} replikacji).}
\end{table}

\subsection{Zmiana na poziomie pracowników}
Spośród \EmpN{} pracowników (\EmpSharePct\% wszystkich), którzy mieli co najmniej jedną sekwencję
liczącą trzy lub więcej prób, \EmpImproved{} poprawiło wynik, \EmpWorsened{} go pogorszyło, a u
\EmpUnchanged{} nie zmienił się on między pierwszą a trzecią próbą (test znaków: $p \SignP$; test
Wilcoxona: $p \WilcoxonP$; korelacja rangowo-dwuseryjna $r = \RankBiserial$). Średnia zmiana wyniosła
\EmpMeanChange{} punktu procentowego (95\% PU \EmpMeanChangeCI; rys.~\ref{fig:zmiana_pracownikow},
tab.~\ref{tab:zmiana_pracownikow}).

\begin{figure}[tbp]
\centering
\includegraphics{\FigDir rys_zmiana_pracownikow}
\caption{Rozkład zmiany odsetka poprawnych odpowiedzi na poziomie pracownika między próbą~3 a
próbą~1 w jego własnych sekwencjach liczących co najmniej trzy próby ($n = \EmpN$ pracowników;
przedziały o szerokości 5 punktów procentowych wyśrodkowane na zerze).}
\label{fig:zmiana_pracownikow}
\end{figure}

\begin{table}[htbp]
\centering
\caption{Zmiana odsetka poprawnych odpowiedzi między pierwszą a trzecią próbą tych samych pytań.}
\label{tab:zmiana_pracownikow}
\small
\tabelacsv{lr}{zmiana_pracownikow}
\uwagatabeli{Zmiana na poziomie pracownika: różnica odsetka poprawnych odpowiedzi między próbą~3
a próbą~1 we własnych sekwencjach pracownika liczących $\geq$3 próby; remisy wykluczono z testu znaków
i testu Wilcoxona; p.p.: punkty procentowe.}
\end{table}

\subsection{Przenoszenie na nowe pytania (PB4)}
Odsetek poprawnych odpowiedzi przy pierwszym kontakcie z nowym pytaniem nie rósł wraz ze stażem
(model~3: OR \ORMThreeTenure{} na rok, 95\% PU \ORMThreeTenureCI; tab.~\ref{tab:regresja}), a po
uwzględnieniu numeru próby staż nie wiązał się dodatnio z poprawnością w próbach powtórnych
(model~2: OR \ORMTwoTenure, 95\% PU \ORMTwoTenureCI). Korzyści wydają się więc specyficzne dla
powtarzanych pytań. % SPRAWDŹ

\subsection{Analiza pytań (PB5)}
Mediana trudności pytań wynosiła \ItemDifficultyMedian\% poprawnych odpowiedzi w pierwszej próbie:
na \ItemsHard{} pytania poprawnie odpowiedziało mniej niż 50\% pracowników, na \ItemsModerate{} od
50\% do 90\%, a na \ItemsEasy{} ponad 90\% (tab.~\ref{tab:pytania}). Mediana korelacji pytanie--reszta
wynosiła \ItemDiscMedian; \ItemsLowDisc{} z \ItemsDiscTested{} pytań miało korelację poniżej 0{,}20,
natomiast w \ItemsGoodDisc{} osiągnęła ona co najmniej 0{,}30 (rys.~\ref{fig:pytania}). Błędy koncentrowały się na jednym
dystraktorze (mediana \TopDistractorMedian\% błędów w pytaniu), a \NNonFunctional{} z \NDistractors{}
zaobserwowanych dystraktorów (\NonFunctionalPct\%) było niefunkcjonalnych. Poprawność istotnie rosła
wraz z powtórkami w \ItemsImproving{} z \ItemsTestable{} testowalnych pytań, a malała w
\ItemsWorsening{} (po korekcie FDR). Najtrudniejsze pytania zestawiono w
tab.~\ref{tab:najtrudniejsze_pytania}.

\begin{figure}[tbp]
\centering
\includegraphics{\FigDir rys_analiza_pytan}
\caption{Trudność pytania (odsetek poprawnych w pierwszej próbie) a dyskryminacja (korelacja
pytanie--reszta). Trójkąty: pytania, w których poprawność istotnie rosła wraz z powtórkami
(Benjamini--Hochberg, $q < \Alpha$). Linie kropkowane: $r = 0{,}20$ i 90\% poprawnych.}
\label{fig:pytania}
\end{figure}

\begin{table}[htbp]
\centering
\caption{Analiza pytań (klasyczna teoria testu).}
\label{tab:pytania}
\small
\tabelacsv{p{8.5cm}r}{pytania}
\uwagatabeli{Trudność i dyskryminację obliczono na pierwszych próbach; korelacja pytanie--reszta
wykorzystuje poprawność pracownika w pozostałych pytaniach (pracownicy z $\geq$10 pierwszymi próbami).
Trend: regresja logistyczna poprawności względem log$_2$(numeru próby) dla każdego pytania z błędami
standardowymi klastrowanymi względem pracowników, korekta Benjaminiego--Hochberga. Dystraktory, których
nigdy nie wybrano, nie są widoczne w logu i nie zostały uwzględnione.}
\end{table}

% \begin{table*}[htbp]
% \centering
% \caption{\TopItems{} najtrudniejszych pytań.}
% \label{tab:najtrudniejsze_pytania}
% \small
% \setlength{\tabcolsep}{3.5pt}
% \tabelacsv{rlrrrrrrr}{najtrudniejsze_pytania}
% \uwagatabeli{Pierwsza/Kolejne: odsetek poprawnych w pierwszej i w późniejszych próbach; Dom.\
% dystr.: udział błędów przypadających na najczęstszy dystraktor; $r_{\mathrm{it}}$: korelacja
% pytanie--reszta; OR trendu: iloraz szans na podwojenie numeru próby; $q$: wartość $p$ skorygowana
% metodą Benjaminiego--Hochberga.}
% \end{table*}

\subsection{Analizy wrażliwości}
Oszacowanie efektu numeru próby pozostawało zbliżone we wszystkich specyfikacjach (OR od \ORRange{};
rys.~\ref{fig:wrazliwosc}, tab.~\ref{tab:wrazliwosc}) i wynosiło \ORSensCohort{} w zbalansowanej kohorcie,
\ORSensLenThree, \ORSensLenFive{} i \ORSensLenEight{} w sekwencjach liczących 3--4, 5--7 i co najmniej
8 prób, \ORSensGee{} w uśrednionym populacyjnie modelu GEE oraz \ORSensMixed{} w modelu mieszanym.
Oszacowanie było wyższe po korekcie o staż (\ORSensTenure), który jest skorelowany z numerem próby.

\begin{figure}[tbp]
\centering
\includegraphics{\FigDir rys_wrazliwosc}
\caption{Iloraz szans (95\% PU) poprawnej odpowiedzi na podwojenie numeru próby w różnych
specyfikacjach modelu. Romb: model główny; linia pionowa: brak efektu.}
\label{fig:wrazliwosc}
\end{figure}

\begin{table}[htbp]
\centering
\caption{Analizy wrażliwości: iloraz szans poprawnej odpowiedzi na podwojenie numeru próby.}
\label{tab:wrazliwosc}
\small
\tabelacsv{p{8cm}rr}{wrazliwosc}
\uwagatabeli{Odchylenia standardowe losowych wyrazów wolnych (skala logitowa): pracownik
SD = \SDEmployee, pytanie SD = \SDItem; ICC w skali latentnej: pracownik = \ICCEmployee, pytanie =
\ICCItem; $R^2$ Nakagawy: brzegowy = \RsqMarginal, warunkowy = \RsqConditional. Warstwy według
długości sekwencji służą kontroli obciążenia typu \emph{pattern-mixture} (ubywania obserwacji)
w krzywej uczenia.}
\end{table}

% =============================================================================
% Dyskusja i wnioski. Plik edytowany ręcznie -- skrypt analiza_odpowiedzi_naukowa.py go NIE nadpisuje.
% Liczby: definicje.tex; literatura: references.bib.
% =============================================================================
\section{Dyskusja}

\subsection{Główne wyniki}
W ciągu trzech lat rutynowego stosowania programu mikrouczenia w miejscu pracy odsetek poprawnych
odpowiedzi na wielokrotnie prezentowane pytania rósł wraz z liczbą przypomnień zgodnie z funkcją
potęgową, a efekt był odporny na ubywanie obserwacji, specyfikację modelu i klastrowanie
danych. Poprawność malała wraz z odstępem od poprzedniego przypomnienia, wcześniejsze błędy były
w kolejnej próbie korygowane tylko częściowo, a powtarzane błędy w większości odtwarzały ten sam
dystraktor. Poprawność odpowiedzi na nowe pytania nie rosła wraz ze stażem. % SPRAWDŹ

\subsection{Porównanie z wcześniejszymi badaniami}
Przebieg przyrostu odpowiada krzywym uczenia opisanym dla wielu zbiorów danych
edukacyjnych \citep{koedinger2023} i jest zgodny z efektami testowania i rozłożenia w czasie
\citep{roediger2006,cepeda2006}. Randomizowane badania edukacji rozłożonej w czasie wykazały lepsze
utrwalenie wiedzy \citep{kerfoot2007mededuc,kerfoot2009}; przedstawione dane obserwacyjne rozszerzają te
ustalenia na niemedyczne środowisko pracy z długimi odstępami między powtórkami. Ujemny związek między
odstępem a poprawnością jest zgodny z modelami krzywej zapominania dopasowywanymi do logów uczących się
\citep{settles2016,tabibian2019}. Podobnie jak w badaniach, w których wielokrotnie wykorzystywano te
same pytania \citep{donker2022}, poprawy w powtarzanych pytaniach nie da się oddzielić od zapamiętania
konkretnego pytania i jego odpowiedzi; brak efektu stażu w nowych pytaniach sugeruje ograniczony
transfer.

\subsection{Implikacje praktyczne}
Program ponownie prezentował pytania w tym samym odstępie niezależnie od wyniku poprzedniej próby.
Harmonogramy adaptacyjne, które skracają odstęp po błędach i wydłużają go po poprawnych odpowiedziach,
efektywniej wykorzystują historię uczącego się \citep{lindsey2014,settles2016,tabibian2019}. Wybory
tych samych dystraktorów utrzymywały się, mimo że po każdej odpowiedzi wyświetlano wyjaśnienie poprawnej odpowiedzi,
co sugeruje, że ogólne wyjaśnienie nie wystarcza do skorygowania błędnych przekonań; skuteczniejsza może
być informacja zwrotna odnosząca się do konkretnego wybranego dystraktora. Analiza pytań wskazała
kandydatów do zmiany: pytania o niskiej dyskryminacji, pytania, na które niemal wszyscy odpowiadali
poprawnie już w pierwszej próbie (mało miejsca na uczenie się), oraz dystraktory niefunkcjonalne
\citep{tarrant2009,gierl2017}.

\subsection{Ograniczenia}
Badanie ma charakter obserwacyjny i nie obejmuje grupy kontrolnej, nie można więc przypisać przyrostu
wiedzy programowi w sensie przyczynowym; poprawa w powtarzanych pytaniach może częściowo odzwierciedlać
rozpoznawanie pytania, a nie jego zrozumienie. Udział i aktywność różniły się między pracownikami,
a wyniki dla sekwencji z wieloma próbami opisują bardziej aktywnych uczestników, choć analizy
zbalansowanej kohorty i długości sekwencji nie wskazały na obciążenie wynikające z ubywania obserwacji.
Odsetek poprawnych odpowiedzi w pierwszych próbach (\AccFirst\%) ogranicza zakres możliwej poprawy
(efekt sufitowy).
Dystraktory, których nigdy nie wybrano, nie są widoczne w logu, więc funkcjonowanie dystraktorów może
być przeszacowane. Nie dysponowano informacjami o stanowiskach pracowników ani ich wcześniejszej wiedzy,
a dane o wieku i płci były niedostępne dla \DemoUnmatched{} z \DemoN{} pracowników; cech
demograficznych nie uwzględniano w modelach. Znaczniki czasu nie zawierały informacji o strefie
czasowej. Model mieszany zawierał wyłącznie losowe wyrazy wolne; nie modelowano indywidualnego tempa uczenia się pracowników.
Dane pochodzą z jednej organizacji.

\section{Wnioski}
Powtarzanemu przypominaniu w ramach rutynowych szkoleń pracowniczych towarzyszyła poprawa ograniczona
do powtarzanych pytań oraz zapominanie między powtórkami. Dane z logów takich programów
mogą służyć do monitorowania uczenia się, oceny pytań i projektowania adaptacyjnych harmonogramów
powtórek.


\section*{Oświadczenia}

% SPRAWDŹ: sposób kodowania danych pracowników oraz status zgody
% komisji etycznej lub zwolnienia z obowiązku jej uzyskania.
\noindent\textbf{Etyka.}
W badaniu wykorzystano dane przekazane przez bank, w których pracownicy
byli reprezentowani przez kody. Dane dotyczące wieku i płci raportowano
wyłącznie w postaci zagregowanej.
\fillin{Informacja o zgodzie właściwej komisji etycznej lub zwolnieniu
z obowiązku jej uzyskania, jeśli jest wymagana.}\\
\textbf{Dostępność danych.}
Dane wykorzystane w badaniu stanowią własność klienta i nie są
udostępniane publicznie.\\
\textbf{Finansowanie.}
Badanie nie otrzymało finansowania.\\
\textbf{Konflikt interesów.}
Autorzy deklarują brak konfliktu interesów.\\
\textbf{Wykorzystanie narzędzi AI.}
Narzędzia generatywnej sztucznej inteligencji wykorzystano do
przygotowania opisu wyników analizy oraz tłumaczenia tekstu
na język angielski.
\bibliography{references}
\end{document}
